% !TEX root =./ECE444_Practice_main.tex
%
% L1 practice set (Course Introduction) -- a short identification/concept set, as
% the practice contract allows for an intro lesson: recognize the six antenna
% families from a geometry description, and state what an antenna is, what it
% does not do, and what reciprocity buys a receiver. No numeric parts, so no
% answer boxes; expected answers are one or two sentences each.

\fullwidth{\textbf{Learning Objective 1.1: } I can explain what an antenna is,
describe its role in a wireless system, and recognize common antenna types by
sight.
\\[4pt]\normalfont\small Show your work. A \textbf{Documentation} line at the end
is required (write \textbf{None} if you did not collaborate).}

\begin{questions}

%%%%%%%%%%%%%%%%%%%%%%%%%%%%%%%%%%%%%%%%%%%%%%%%%%%%%%%%%%%%%%%%%%%%%%%%%%%%
\question \textbf{Name that antenna.} Each description below is what you would
see if you walked up to the hardware. For each, name the antenna family and give
its \emph{job} --- the reason an engineer reaches for that shape --- in one or two
sentences.

\begin{parts}

	\part \textbf{(i)} A straight metal rod, split and fed at its centre, roughly
	half a wavelength long overall. \textbf{(ii)} A single rod about a quarter of a
	wavelength tall, standing on a large sheet of metal or a vehicle roof and fed at
	its base.
		\begin{solution}[1.6in]
		\textbf{(i) Half-wave dipole.} The reference element of the whole subject: it
		is simple, resonant, easy to feed, and radiates in a doughnut around the wire,
		which means it covers every azimuth direction at once. Nearly every gain figure
		you will meet is quoted against it or against the isotropic radiator.
		\textbf{(ii) Quarter-wave monopole.} The same antenna with the bottom half
		replaced by an image in the ground plane, so it does the dipole's job at half
		the height. That is why it is the default on vehicles, aircraft skins, and
		handheld radios, where physical height is the scarce resource.
		\end{solution}

	\part A thin rectangle of copper etched on a grounded dielectric board, about
	half a wavelength across, fed from underneath or from a trace at its edge.
		\begin{solution}[1.4in]
		A \textbf{microstrip patch}. Its job is to be \emph{flat}: it conforms to a
		skin or a handset case, costs nothing to reproduce once the board is designed,
		and tiles neatly into arrays. The price is narrow bandwidth and a pattern that
		only covers the hemisphere above the board --- which is exactly what a GPS
		receiver or a phased-array element wants anyway.
		\end{solution}

	\part A rectangular metal pipe that flares open at one end, fed from a waveguide
	at the other.
		\begin{solution}[1.4in]
		A \textbf{horn}. Its job is to be trustworthy: a horn is a controlled,
		well-understood aperture with moderate gain, low loss, and a clean pattern over
		a wide band, so it serves as the standard-gain reference on a test range and as
		the feed that illuminates a reflector. It is bulky, which is why you rarely see
		one bolted to an aircraft.
		\end{solution}

	\part \textbf{(i)} A row of parallel rods of slightly different lengths mounted
	along a boom, with only one of them connected to a cable. \textbf{(ii)} A large
	curved metal dish with a small antenna suspended at its focus.
		\begin{solution}[1.7in]
		\textbf{(i) Yagi-Uda array.} Its job is cheap directional gain: the unfed
		parasitic rods re-radiate with the right phase to reinforce the pattern forward
		and cancel it backward, so you get a real beam from a single feed point and a
		bag of aluminium. \textbf{(ii) Reflector (dish) antenna.} Its job is the
		highest gain per dollar at microwave frequencies: the small antenna at the
		focus is only the \emph{feed}, and the dish is the aperture that does the work.
		Gain grows with the dish area, so you buy beam by buying diameter.
		\end{solution}

\end{parts}

\newpage
%%%%%%%%%%%%%%%%%%%%%%%%%%%%%%%%%%%%%%%%%%%%%%%%%%%%%%%%%%%%%%%%%%%%%%%%%%%%
\question \textbf{What an antenna is, and what it is for.} Answer each part in
one or two complete sentences.

\begin{parts}

	\part Define an antenna as a \emph{transducer}. Name the two kinds of wave it
	sits between, and say which component in an RF link block diagram
	(source $\rightarrow$ amplifier $\rightarrow$ transmission line
	$\rightarrow$ antenna $\rightarrow$ channel) is the only one that touches free
	space.
		\begin{solution}[1.4in]
		An antenna is the transducer between a \textbf{guided wave} --- a signal
		travelling on a cable or in a waveguide, confined by conductors --- and a
		\textbf{radiating wave} travelling freely through space. It is the only
		component in the chain that touches free space; everything upstream of it is
		guided-wave hardware, which is why an antenna problem cannot be repaired
		anywhere else in the radio.
		\end{solution}

	\part Antennas do not create energy. Given that, explain what a vendor means by
	``this antenna has 30 dBi of gain.''
		\begin{solution}[1.4in]
		Gain is a statement about \emph{direction}, not about power. The antenna
		radiates whatever the amplifier hands it (less its own losses); $30\dbi$ says
		that in its best direction the power density is $1000$ times what an isotropic
		radiator fed the same power would produce there. That concentration is paid for
		everywhere else --- a high-gain antenna radiates far \emph{less} than isotropic
		over most of the sphere.
		\end{solution}

	\part A tracking antenna is measured on transmit and found to have a $3\mydeg$
	beam pointed at boresight. You now use the same antenna to \emph{receive}. State
	what reciprocity tells you about the receive pattern, and give two practical
	consequences for the receiving station.
		\begin{solution}[1.6in]
		Reciprocity says the receive pattern is the \emph{same} pattern: the antenna is
		$3\mydeg$ wide on receive too, with the same gain and the same sidelobes. Two
		consequences. \textbf{Pointing:} the station must aim within a fraction of a
		degree or the signal falls off the beam --- receive sensitivity is directional
		in exactly the way transmit power is. \textbf{Interference and noise:} the same
		narrow beam that concentrates transmitted power also rejects signals and thermal
		noise arriving from every other direction, so a high-gain antenna improves the
		link twice over --- once at each end.
		\end{solution}

	\part Match each band to the mission it is usually flown for, and state in one
	sentence what sets the size of the antenna in each case.
	\textbf{Bands:} (i) $3$--$30\mhz$ (HF), (ii) $225$--$400\mhz$ (UHF),
	(iii) $1575\mhz$ (GPS L1), (iv) $8$--$12\ghz$ (X-band).
	\textbf{Missions:} fire-control and imaging radar; beyond-line-of-sight skywave
	communications; navigation; tactical manpack radio and UHF SATCOM.
		\begin{solution}[1.8in]
		\textbf{(i) HF $\rightarrow$ beyond-line-of-sight skywave communications},
		bouncing off the ionosphere when no satellite is available.
		\textbf{(ii) UHF $\rightarrow$ tactical manpack radio and UHF SATCOM}, the band
		where a whip is short enough to carry but the wavelength still penetrates
		foliage and buildings.
		\textbf{(iii) $1575\mhz$ $\rightarrow$ navigation} (GPS L1), served by a small
		circularly polarized patch.
		\textbf{(iv) X-band $\rightarrow$ fire-control and imaging radar}, where a
		modest dish still gives a very narrow beam.
		What sets the size in every case is the \textbf{wavelength}: a resonant element
		is a fixed fraction of $\lambda$, and an aperture needs many $\lambda$ across to
		make a narrow beam. An HF antenna is tens of metres and a GPS patch is a few
		centimetres for exactly the same reason.
		\end{solution}

\end{parts}

\vspace{1.5em}
\noindent\textbf{Documentation:} \rule[-2pt]{0.6\linewidth}{0.4pt}

\end{questions}
